\documentclass[11pt]{article}
\usepackage[margin=2.5cm]{geometry}
\usepackage{amsmath, amssymb, amsthm}
\usepackage{bm}
\usepackage{hyperref}

\newcommand{\xb}{\bm{x}}
\newcommand{\eps}{\bm{\varepsilon}}
\newcommand{\abar}{\bar{\alpha}}
\newcommand{\E}{\mathbb{E}}
\newcommand{\R}{\mathbb{R}}
\newcommand{\N}{\mathcal{N}}

\title{Denoising Diffusion Probabilistic Models (DDPMs)}
\author{Notes for ChemDM}
\date{}

\begin{document}
\maketitle

\section{Setup}

We have data $\xb_0 \sim q(\xb_0)$ in $\R^d$ and want to learn a generative
model that can produce new samples from $q(\xb_0)$.  The idea is to
\emph{gradually destroy} the data by adding noise (the forward process) and
then \emph{learn to reverse} that destruction (the reverse process).

\section{Forward process}

Define a sequence of noise levels through a \textbf{variance schedule}
$\beta_1, \beta_2, \ldots, \beta_T \in (0,1)$.  The forward process adds
Gaussian noise in $T$ steps:
\begin{equation}
    q(\xb_t \mid \xb_{t-1})
    = \N\!\left(\xb_t;\; \sqrt{1 - \beta_t}\,\xb_{t-1},\; \beta_t \bm{I}\right).
\end{equation}
Each step shrinks the signal by $\sqrt{1-\beta_t}$ and adds noise with
variance $\beta_t$.  After many steps, $\xb_T$ is approximately standard
Gaussian regardless of $\xb_0$.

\subsection{Closed-form sampling at arbitrary $t$}

Define $\alpha_t = 1 - \beta_t$ and $\abar_t = \prod_{s=1}^{t} \alpha_s$.
Then we can jump directly from $\xb_0$ to $\xb_t$ without simulating
intermediate steps:
\begin{equation}\label{eq:q_sample}
    q(\xb_t \mid \xb_0)
    = \N\!\left(\xb_t;\; \sqrt{\abar_t}\,\xb_0,\;
      (1 - \abar_t)\bm{I}\right),
\end{equation}
or equivalently,
\begin{equation}\label{eq:reparam}
    \xb_t = \sqrt{\abar_t}\,\xb_0 + \sqrt{1 - \abar_t}\,\eps,
    \qquad \eps \sim \N(\bm{0}, \bm{I}).
\end{equation}
This is the \textbf{reparameterization trick} and is what we use during
training to generate noisy samples.

\paragraph{Proof sketch.}
Each forward step is a linear Gaussian map.  Composing $t$ of them gives
$\xb_t = \bigl(\prod_{s=1}^t \sqrt{\alpha_s}\bigr)\xb_0 + \text{noise}$,
with the noise terms accumulating into a single Gaussian with variance
$1 - \prod_{s=1}^t \alpha_s = 1 - \abar_t$.

\subsection{Variance schedule}

Two common choices:
\begin{itemize}
    \item \textbf{Linear:} $\beta_t$ increases linearly from $\beta_1
          \approx 10^{-4}$ to $\beta_T \approx 0.02$.
    \item \textbf{Cosine} (Nichol \& Dhariwal, 2021): define $\abar_t$
          directly as
          \[
              \abar_t = \frac{f(t)}{f(0)},
              \qquad f(t) = \cos^2\!\left(\frac{t/T + s}{1+s}
              \cdot \frac{\pi}{2}\right),
          \]
          with a small offset $s=0.008$ to avoid $\beta_t$ becoming too
          small near $t=0$.  Then $\beta_t = 1 - \abar_t / \abar_{t-1}$.
          This schedule spends more of the diffusion budget at intermediate
          noise levels, where the denoising task is most informative.
\end{itemize}

\section{Reverse process}

The key insight: if $\beta_t$ is small, each forward step
$q(\xb_t \mid \xb_{t-1})$ is nearly Gaussian, so the true reverse step
$q(\xb_{t-1} \mid \xb_t)$ is \emph{also} nearly Gaussian.  We can therefore
parameterize a learned reverse:
\begin{equation}
    p_\theta(\xb_{t-1} \mid \xb_t)
    = \N\!\left(\xb_{t-1};\; \bm{\mu}_\theta(\xb_t, t),\;
      \sigma_t^2 \bm{I}\right),
\end{equation}
where $\bm{\mu}_\theta$ is a neural network and $\sigma_t^2$ is set to
a known value (see below).  Sampling starts from $\xb_T \sim \N(\bm{0},
\bm{I})$ and iterates $\xb_{t-1} \sim p_\theta(\xb_{t-1} \mid \xb_t)$ for
$t = T, T{-}1, \ldots, 1$.

\subsection{The true posterior}

If we \emph{knew} $\xb_0$, the reverse step would be tractable:
\begin{equation}\label{eq:posterior}
    q(\xb_{t-1} \mid \xb_t, \xb_0)
    = \N\!\left(\xb_{t-1};\; \tilde{\bm{\mu}}_t(\xb_t, \xb_0),\;
      \tilde{\beta}_t \bm{I}\right),
\end{equation}
with
\begin{align}
    \tilde{\bm{\mu}}_t(\xb_t, \xb_0)
    &= \frac{\sqrt{\abar_{t-1}}\,\beta_t}{1 - \abar_t}\,\xb_0
     + \frac{\sqrt{\alpha_t}\,(1 - \abar_{t-1})}{1 - \abar_t}\,\xb_t,
    \label{eq:posterior_mean} \\[6pt]
    \tilde{\beta}_t
    &= \frac{(1 - \abar_{t-1})}{1 - \abar_t}\,\beta_t.
    \label{eq:posterior_var}
\end{align}
The posterior mean is a weighted average of $\xb_0$ (the clean data) and
$\xb_t$ (the current noisy state).

\section{Training objective}

\subsection{Variational bound}

The DDPM training objective is derived from the variational lower bound
(ELBO) on $\log p_\theta(\xb_0)$.  After standard manipulations, it
decomposes into $T$ KL-divergence terms, each comparing the learned reverse
$p_\theta(\xb_{t-1} \mid \xb_t)$ against the true posterior
$q(\xb_{t-1} \mid \xb_t, \xb_0)$.  Since both are Gaussians with the
same (fixed) variance $\tilde\beta_t$, the KL reduces to a squared
difference between their means.

The upshot: training reduces to teaching $\bm\mu_\theta(\xb_t, t)$ to
approximate $\tilde{\bm\mu}_t(\xb_t, \xb_0)$.

\subsection{Parameterization choices}

We can parameterize $\bm\mu_\theta$ in several equivalent ways.  From~\eqref{eq:reparam} we can express $\xb_0$ and $\eps$ in terms of each
other:
\begin{equation}
    \xb_0 = \frac{\xb_t - \sqrt{1-\abar_t}\,\eps}{\sqrt{\abar_t}},
    \qquad
    \eps = \frac{\xb_t - \sqrt{\abar_t}\,\xb_0}{\sqrt{1-\abar_t}}.
\end{equation}

\paragraph{$\eps$-prediction (Ho et al., 2020).}
Let the network predict the noise: $\eps_\theta(\xb_t, t) \approx \eps$.
Substituting into \eqref{eq:posterior_mean}:
\begin{equation}
    \bm{\mu}_\theta(\xb_t, t)
    = \frac{1}{\sqrt{\alpha_t}}
      \left(\xb_t - \frac{\beta_t}{\sqrt{1-\abar_t}}\,
      \eps_\theta(\xb_t, t)\right).
\end{equation}
The simplified training loss is:
\begin{equation}\label{eq:eps_loss}
    L_\text{simple}
    = \E_{t,\, \xb_0,\, \eps}\!\left[
        \bigl\|\eps - \eps_\theta(\xb_t, t)\bigr\|^2
      \right],
\end{equation}
where $t \sim \mathrm{Uniform}\{1,\ldots,T\}$, $\xb_0 \sim q(\xb_0)$,
$\eps \sim \N(\bm{0},\bm{I})$, and $\xb_t$ is given
by~\eqref{eq:reparam}.

\paragraph{$\xb_0$-prediction.}
Let the network predict the clean data directly:
$\hat{\xb}_0 = f_\theta(\xb_t, t) \approx \xb_0$.  The training loss is:
\begin{equation}\label{eq:x0_loss}
    L_{x_0}
    = \E_{t,\, \xb_0,\, \eps}\!\left[
        \bigl\|\xb_0 - f_\theta(\xb_t, t)\bigr\|^2
      \right].
\end{equation}
At sampling time, plug the predicted $\hat\xb_0$ into the posterior
mean~\eqref{eq:posterior_mean} to get $\bm\mu_\theta$, then sample:
\begin{equation}
    \xb_{t-1}
    = \tilde{\bm\mu}_t(\xb_t,\, \hat\xb_0)
    + \sqrt{\tilde\beta_t}\,\bm{z},
    \qquad \bm{z} \sim \N(\bm{0}, \bm{I}),
\end{equation}
with no noise added at $t=1$ (i.e.\ we return $\hat\xb_0$ directly).

\paragraph{Equivalence.}
The two parameterizations are related by the linear
map~\eqref{eq:reparam}.  The losses~\eqref{eq:eps_loss}
and~\eqref{eq:x0_loss} differ only by a $t$-dependent reweighting factor
$\abar_t / (1 - \abar_t)$.  In practice, $\xb_0$-prediction tends to work
better when combined with a cosine schedule, while $\eps$-prediction is the
original DDPM formulation.

\section{Training algorithm}

The full training procedure is:
\begin{enumerate}
    \item Sample a data point $\xb_0 \sim q(\xb_0)$.
    \item Sample a random timestep $t \sim \mathrm{Uniform}\{0, 1,
          \ldots, T{-}1\}$.
    \item Sample noise $\eps \sim \N(\bm{0}, \bm{I})$.
    \item Compute $\xb_t = \sqrt{\abar_t}\,\xb_0
          + \sqrt{1-\abar_t}\,\eps$.
    \item Predict $\hat{\xb}_0 = f_\theta(\xb_t, t)$
          \quad (or $\hat{\eps} = \eps_\theta(\xb_t, t)$).
    \item Compute loss: $\|\xb_0 - \hat{\xb}_0\|^2$
          \quad (or $\|\eps - \hat{\eps}\|^2$).
    \item Backpropagate and update $\theta$.
\end{enumerate}
Note that each gradient step uses a \emph{single random} $t$.  Over many
iterations the network learns to denoise at all noise levels.

\section{Sampling algorithm}

\begin{enumerate}
    \item Sample $\xb_T \sim \N(\bm{0}, \bm{I})$.
    \item For $t = T, T{-}1, \ldots, 1$:
    \begin{enumerate}
        \item Predict $\hat\xb_0 = f_\theta(\xb_t, t)$.
        \item Compute the posterior mean $\tilde{\bm\mu}_t$ using
              \eqref{eq:posterior_mean} with $\hat\xb_0$.
        \item Sample $\bm{z} \sim \N(\bm{0}, \bm{I})$ if $t > 1$, else
              $\bm{z} = \bm{0}$.
        \item Set $\xb_{t-1} = \tilde{\bm\mu}_t
              + \sqrt{\tilde\beta_t}\,\bm{z}$.
    \end{enumerate}
    \item Return $\xb_0$.
\end{enumerate}

\section{Connection to score-based models}

The \textbf{score} of a distribution is $\nabla_{\xb} \log p(\xb)$.  For
the forward marginal $q(\xb_t \mid \xb_0)$ from~\eqref{eq:q_sample}:
\begin{equation}
    \nabla_{\xb_t} \log q(\xb_t \mid \xb_0)
    = -\frac{\xb_t - \sqrt{\abar_t}\,\xb_0}{1 - \abar_t}
    = -\frac{\eps}{\sqrt{1 - \abar_t}}.
\end{equation}
So predicting $\eps$ is equivalent to estimating the score (up to a known
scaling).  Score-based generative models (Song \& Ermon, 2019) take this
perspective and work in continuous time, replacing the discrete chain with
a stochastic differential equation (SDE).  The DDPM discrete chain can be
seen as an Euler--Maruyama discretization of that SDE.

\section{Application to transition paths}

In our setting, $\xb_0 = \xb(s) \in \R^{N_\text{atoms} \times 3}$ is the
atomic configuration at normalized arclength $s$ along a transition path.
The denoising network $f_\theta$ is conditioned on:
\begin{itemize}
    \item $\xb_A$, $\xb_B$: reactant and product geometries (fixed graphs),
    \item $s \in [0,1]$: position along the path,
    \item $t$: diffusion noise level.
\end{itemize}
The network predicts the clean configuration $\hat\xb_0 = f_\theta(\xb_t,
\xb_A, \xb_B, s, t)$.

Because the reverse process is \emph{stochastic}, sampling with the same
$(\xb_A, \xb_B, s)$ but different random seeds produces
\textbf{different transition paths}.  This is precisely the property we
want: for a given reaction, multiple mechanistically valid pathways can
exist, and a diffusion model can represent this multi-modality naturally,
whereas a deterministic regression model is forced to average over them.

At the boundaries, the conditioning enforces $\hat\xb_0 \approx \xb_A$
when $s \approx 0$ and $\hat\xb_0 \approx \xb_B$ when $s \approx 1$,
regardless of the noise level $t$.

\section*{Key references}
\begin{itemize}
    \item Ho, Jain \& Abbeel. \textit{Denoising Diffusion Probabilistic
          Models.} NeurIPS, 2020.
    \item Nichol \& Dhariwal. \textit{Improved Denoising Diffusion
          Probabilistic Models.} ICML, 2021.
    \item Song \& Ermon. \textit{Generative Modeling by Estimating Gradients
          of the Data Distribution.} NeurIPS, 2019.
    \item Song et al. \textit{Score-Based Generative Modeling through
          Stochastic Differential Equations.} ICLR, 2021.
\end{itemize}

\end{document}
